\documentclass[11pt]{article}

%------------------------------------------------
% Packages
%------------------------------------------------
\usepackage[a4paper,margin=2.5cm]{geometry}
\usepackage{amsmath,amssymb,amsthm}
\usepackage{mathtools}
\usepackage{enumitem}
\usepackage{hyperref}

\hypersetup{
    colorlinks=true,
    linkcolor=blue,
    citecolor=blue,
    urlcolor=blue
}

%------------------------------------------------
% Title Information
%------------------------------------------------

\title{\bfseries Point Source Localisation with Unbalanced Optimal Transport}

\author{
Tuomo Valkonen\thanks{MODEMAT
\texttt{tuomo.valkonen@iki.fi}}
\and
Thi Tam Dang\thanks{Department of Mathematics and Statistics University of Helsinki,
\texttt{tam.dang@iki.fi}}
}

\date{}

%------------------------------------------------

\begin{document}

\maketitle

\vspace{-1em}

\noindent
\textbf{Keywords:}
optimization, inverse problems, measure spaces, convectionn-diffusion

\vspace{0.5em}


\vspace{1em}

%%%%%%%%%%%%%%%%%%%%%%%%%%%%%%%%%%%%%%%%%%%%%%%%%%
\section*{Abstract}

Replacing the quadratic proximal penalty familiar from Hilbert spaces by an unbalanced optimal transport distance, we develop forward-backward type optimisation methods in spaces of Radon measures. We avoid the actual computation of the optimal transport distances through the use of transport three-plans and the rough concept of transport subdifferentials. The resulting algorithm has a step similar to the sliding heuristics previously introduced for conditional gradient methods, however, now non-heuristically derived from the geometry of the space \cite{tuomov-unbalanced}.

Afterwards, as an application of the above algorithm, we study the inverse problem of locating gas leaks from line-of-sight concentration measurements using a convection–diffusion model with the source term a Radon measure. By imposing sparsity-promoting regularisation on this measure, we recover point sources---identifying both their locations and intensities---rather than diffuse approximations.
We jointly estimate the underlying physical convection (wind) and diffusion parameters.
Our main theoretical contribution is the stability analysis of the convection--diffusion equation with respect to its parameters: the measure, and the convection and diffusion fields.
Our experiments demonstrate accurate localisation, highlighting the potential of the method for practical gas emission detection \cite{dangvalkonen2026leak}.


%%%%%%%%%%%%%%%%%%%%%%%%%%%%%%%%%%%%%%%%%%%%%%%%%%
\vspace{1em}

\noindent
\textbf{Acknowledgements.}
This work was supported by the Jane and Aatos Erkko Foundation.

%%%%%%%%%%%%%%%%%%%%%%%%%%%%%%%%%%%%%%%%%%%%%%%%%%
\vspace{1em}

\providecommand{\eprint}[1]{\href{http://arxiv.org/abs/#1}{arXiv:#1}}
\begin{thebibliography}{1}

\bibitem{dangvalkonen2026leak}
T.\,T{.\nobreak\kern 0.33333em}Dang and T{.\nobreak\kern 0.33333em}Valkonen,
  Leak localisation with a measure source convection--diffusion model, 2026,
  \href{https://arxiv.org/abs/2605.12095}{\nolinkurl{arXiv:2605.12095}}.
\newblock submitted.

\bibitem{tuomov-unbalanced}
T{.\nobreak\kern 0.33333em}Valkonen, Point source localisation with unbalanced
  optimal transport, 2025,
  \href{https://arxiv.org/abs/2502.12417}{\nolinkurl{arXiv:2502.12417}}.
\newblock submitted.

\end{thebibliography}

\end{document}
